\documentclass[11pt]{article}
\usepackage{fontspec}
\setmainfont{Times New Roman}
\setsansfont{Arial}
\setmonofont{Consolas}
\usepackage{amsmath,amssymb,mathtools}
\usepackage{geometry}
\usepackage{microtype}
\geometry{a4paper,margin=2.28cm}
\setlength{\parindent}{1.15em}
\setlength{\parskip}{0pt}
\makeatletter
\renewcommand\footnotesize{\@setfontsize\footnotesize{7.5}{8.5}\selectfont}
\makeatother
\providecommand{\ma}{\mathfrak{a}}
\providecommand{\mb}{\mathfrak{b}}
\providecommand{\mc}{\mathfrak{c}}
\providecommand{\md}{\mathfrak{d}}
\providecommand{\mC}{\mathfrak{C}}
\providecommand{\mZ}{\mathfrak{Z}}
\providecommand{\mo}{\mathfrak{o}}
\providecommand{\mr}{\mathfrak{r}}
\emergencystretch=120em
\allowdisplaybreaks
\sloppy
\hbadness=10000
\vbadness=10000
\hfuzz=2pt
\vfuzz=10pt
\raggedbottom
\begin{document}
% Unit: Paper34 section12 v001. Direct Interslavic Latin translation from audited German source lines 732--770, segments P34-S0137--P34-S0149. Source scan pages 22--23 retained as witnesses; Section13 heading P34-S0150 is intentionally reserved for Section13. Zenodo record 20822156 was rechecked for this unit at 2026-06-24T22:37Z UTC with no file-list change.
\section*{34. Hiperkompleksne veličiny i teorija predstavjenj}
\emph{Prodolženje: glava II, §§8--14.}

\section*{II. poglavje. Nekomutativna teorija idealov}

\subsection*{§ 12. Nilpotentne idealy}

Ideal $\mc$ nazyva se \emph{nilpotentnym}, ako $\mc^\rho=0$.

\emph{Priměr.} Ideal $(u)$ v §10.

\emph{Ako v $\mo$ jest nilpotentny pravy ideal $\mr$, togda suščestvuje takože nilpotentny lěvy ideal (navět dvustranny ideal).}

\emph{Dokaz.} Nehaj $\mr^\rho=0$. Togda $\mo\mr$ jest lěvy ideal; odnosno, ako $\mo$ ne sadrži jediničnogo elementa i $\mo\mr$ moglo by byti nulovym, lěvym idealom jest $(\mo\mr,\mr)$, i
\[
        (\mo\mr)^\rho=\mo\mr\mo\mr\cdots\mo\mr
        \subseteq\mo\mr\mr\cdots\mr=\mo\mr^\rho=0.
\]

\emph{Suma dvoh nilpotentnyh pravyh idealov jest znovu nilpotentny pravy ideal.}

\emph{Dokaz.} Nehaj $\mc^\rho=0$, $\md^\sigma=0$. V
\[
        (\mc,\md)^{\rho+\sigma-1}=(\ldots,\md\mc\cdots\md\cdots\mc,\ldots)
\]
vsaky člen sadrži ili najmanje $\rho$ faktorov $\mc$, ili najmanje $\sigma$ faktorov $\md$. V prvom slučaju imajemo
\[
        \md\mc\cdots\md\cdots\mc\cdots\subseteq\md\mc\mc\cdots\mc=\md\mc^\rho=0;
\]
v drugom slučaju odpovědajuče za faktory $\md$. Zato
\[
        (\mc,\md)^{\rho+\sigma-1}=0.
\]

Ako teper v $\mo$ važi uslovje maksimalnosti za prave idealy, togda suščestvuje maksimalny nilpotentny pravy ideal. On sadrži vse druge nilpotentne prave idealy, bo inače možno bylo by tvoriti sumu s takim idealom. Nehaj on bude $\mc$. Poneže $\mo\mc$ jest takože nilpotentny pravy ideal, imajemo $\mo\mc\subseteq\mc$, zato $\mc$ jest lěvy ideal. Vsaky nilpotentny lěvy ideal $\md$ takože sadrži se v $\mc$. Bo $(\md,\md\mo)$ jest nilpotentny pravy ideal, zato $\md\subseteq\mc$. \emph{Tako suščestvuje maksimalny (dvustranny) nilpotentny ideal $\mc$, abo “radikal”, ktory sadrži vse druge (prave i lěve).}

V priměru §10 ideal $(u)$ jest maksimalny nilpotentny ideal.

Ako radikal jest nulovy ideal, govori se o “kolcu bez radikala” (“poluprostom kolcu”, “Dedekindovom sistemu”). Kolco klasov ostatkov po radikalu jest vsegda kolco bez radikala.

\emph{Ako $\mZ$ jest centar od $\mo$, a $\mc$ radikal od $\mo$, togda $\mC=\mc\cap\mZ$ jest radikal od $\mZ$.}

\emph{Dokaz.} Jasno jest, že $\mC$ jest nilpotentny. Ako by v $\mZ$ suščestvoval ještě širši nilpotentny ideal $\overline{\mC}$, togda jego možno bylo by razširiti do nilpotentnogo ideala $\overline{\mc}$, ne cělkom sadrženogo v $\mc$, bo
\[
        (\overline{\mC}\mo)^\rho=\overline{\mC}^{\rho}\mo^\rho=0.
\]

\emph{Poslědok.} Ako $\mo$ jest kolco bez radikala, togda takože $\mZ$. Obratno tvrđenje však ne važi, kako pokazuje naš priměr v §10. Centar $\mo/\mc$ sadrži kolco izomorfno k $\mZ/\mC$; ono može však byti vlastnym podkolcom (priměr v §10).

\end{document}
